\documentclass[10pt,twocolumn]{article}

% ------------------ PAQUETES ------------------
\usepackage[spanish]{babel}
\usepackage[utf8]{inputenc}
\usepackage[T1]{fontenc}
\usepackage{geometry}
\usepackage{graphicx}
\usepackage{amsmath, amssymb}
\usepackage{siunitx}
\usepackage{authblk}
\usepackage{abstract}
\usepackage{caption}
\usepackage{titlesec}
\usepackage{multicol}
\usepackage{hyperref}
\usepackage{currfile}
\usepackage{float}

\geometry{margin=1.8cm}

% ------------------ FORMATO DE SECCIONES ------------------
\titleformat{\section}{\large\bfseries}{\thesection}{1em}{}
\titleformat{\subsection}{\normalsize\bfseries}{\thesubsection}{1em}{}

% ------------------ TÍTULO ------------------
\title{\textbf{Aproximación a las estrellas variables}\\Trabajo con métodos de machine learning para análisis de periodicidad\\Segundo trabajo práctico de la serie}


% ------------------ AUTORES ------------------
\author[1]{Bruno Abello Laurent}

\affil[1]{Universidad de Los Andes, Departamento de Física, Colombia}

% Correos (puedes modificar)
\renewcommand\Authands{, }
\date{}

\begin{document}

\twocolumn[
\maketitle
\begin{center}
\small
\textit{Correo institucional:} \href{mailto:{b.abello@uniandes.edu.co}}{b.abello@uniandes.edu.co} \\
\textit{Repositorio de GitHub: } \href{https://github.com/Brunelleschk1/Personal-projects/tree/main/Proyectos\%20de\%20Astronom\%C3\%ADa/Laboratorio\%208}{Repositorio del laboratorio}
\end{center}

\begin{abstract}
\noindent
El presente informe es el tercero del autor con respecto a las estrellas variables. Se centra en el análisis de dos tipos de estrellas, 
las estrellas variables de largo periodo tipo Mira y las estrellas irregulares T Tauri. El propósito de este documento es explicar 
los métodos utilizados por el autor para encontrar los periodos de las estrellas MIRA, y para encontrar la frecuencia envolvente de 
la T Tauri. Se procedió seleccionando una cantidad pequeña de estrellas Mira, unas 5 estrellas del catálogo de OGLE, y trabajando 
con los resultados del catálogo y los propios del autor. De igual manera, de los resultados conocidos de la estrella T Tauri(Su periodo
principal), se procedió a encontrar el periodo de bombeo o pulso de la misma. 
\end{abstract}

\noindent\textbf{Palabras clave:}
\vspace{0.5cm}
]

% ------------------ CONTENIDO ------------------

\section{Introducción}

Este trabajo es sumamente breve en esencia: hablar de 6 estrellas es, en esencia, sencillo, por lo cual este es un documento
más centrado a la discusión de métodos computacionales que de estrellas como tal. 

\subsection{Objetivos}

Los objetivos que tenemos, en términos simples, son los siguientes:
\begin{itemize}
    \item Implementar un método conocido para encontrar el periodo de las estrellas MIRA
    \item Implementar un método desconocido para encontrar el periodo de las estrellas MIRA
    \item Corroborar los datos mostrados en la página de OGLE para las estrellas MIRA seleccionadas
    \item Corroborar, con métodos conocidos, los datos entregados al autor sobre la estrella T Tauri
    \item Implementar maneras de encontrar el periodo de pulsación de la estrella T Tauri
\end{itemize}

\section{Procedimiento}

El proceder para este laboratorio, igual que el anterior, se realizó por partes y sobre una extensión de tiempo considerable. 
Se trabajó en el esqueleto de ambas implementaciones, y más adelante se regresó a completar lo que hacía falta. Contrariamente a
los casos anteriores, este laboratorio conoció un proceder mucho más explorativo: una vez aplicada una solución, se notaban fallos 
y se buscaba una mejor aplicación. Así, se realizaron diversas implementaciones, cada una con un grado de éxito variado. 

\section{Problemas presentados y soluciones aplicadas}

\begin{enumerate}
    \item 
\end{enumerate}


\section{Resultados}



\subsection{Lomb-Scargle}

(Fig: ~\ref{STAR29OGLE}) 

\begin{figure}[h]
    \centering
    \includegraphics[width=1\linewidth]{../OGLE_Results/Star_29.png}
    \caption{Estrella 29 de las analizadas en OGLE. Se nota ruido alrededor de la frecuencia real de la estrella, y se ven 
    caídas en la curva de luz que no parecen ser consistentes sobre todo el periodo de observación de la estrella.}
    \label{STAR29OGLE}
\end{figure}

\section{Conclusiones}


\section*{Comentarios al margen}


\end{document}